\section{Introduction}
\label{sec:intro}

End-to-end autonomous driving is trained on large trajectory corpora under
objectives that favor a single consensus behavior. The resulting policy acts as
one \emph{average driver}: the same deterministic driving style for every user
and scene. Two failures follow at the level of measurable outcomes. First,
planners with similar closed-loop safety scores can occupy very different
regions of a multi-dimensional style profile---speed, acceleration, jerk,
headway, and spacing---on the same scenarios (\cref{fig:teaser}). Optimizing
safety alone therefore collapses diverse, user-relevant ways of driving into one
interchangeable ``good enough'' bucket. Second, when several maneuvers are
equally plausible, learning and decoding that average behavior tend to produce
hesitant, between-mode trajectories instead of a clear commitment. Driving style
is the dimension this paradigm eradicates: not a vague label, but a measurable,
interpretable aspect of how motion is executed, which should vary with user
preference while still respecting the scene.

Much prior work on personalized or styled driving reports joint drops in
closed-loop scores when stylistic behavior is emphasized. That pattern is
widespread, but it does not by itself prove an inevitable trade-off between
style and safety---especially when safety metrics already fail to distinguish
behaviors that users would treat as different. What is missing is a formulation
where style is an explicit, controllable axis rather than an afterthought
entangled with scene understanding and a single-mode decoder.

We aim to measure driving style on shared benchmarks, expose it as a
user-steerable axis, and generate trajectories with a style-conditioned
flow-matching (or diffusion-class) planner so that multi-modal scenes admit
decisive, sample-based choices instead of token-level averaging.

\paragraph{Contributions.}
\begin{itemize}
    \item We provide a multi-planner, human-anchored empirical map of driving
    style and show that a typical end-to-end policy covers only a point on that
    map.
    \item We analyze the costs of the average-driver view, including
    multi-modal compromise.
    \item We present an initial axis-conditioned generative planning direction
    built on a vision--language--action backbone, as a path toward personalized
    driving without treating safety scores as a sufficient description of
    behavior.
\end{itemize}

% Teaser placeholder: style map with iso-safety, different style regions.
\begin{figure}[t]
    \centering
    \fbox{\parbox[c][3.2cm][c]{0.9\linewidth}{\centering\small
    \TODO{Figure 1: human-anchored style map. Planners with similar
    closed-loop safety occupy different style regions; a typical E2E policy is a
    single point.}}}
    \caption{Planners with similar closed-loop safety scores occupy very
    different regions of a multi-dimensional style profile.}
    \label{fig:teaser}
\end{figure}
